% Study sheet: Volume II, Chapter 1, Introduction.
% Numbers: \val{...} from introduction_studysheet_numbers.tex, written by introduction_studysheet_numbers.py.
\documentclass{studysheet}
\input{introduction_studysheet_numbers}
\sheetvolume{2}
\sheetchapter{1}{Introduction}
\sheetsources{Numbers from the mlsysim registry and the chapter's modeled scenarios; questions marked assumed use new inputs.}

\begin{document}
% ===================================================================== FRONT: LEARN
\frontside
\begin{bigidea}
Once a model needs many machines, the \hl{fleet} is the computer: accelerators, network, power and
recovery working as one. Scaling laws say how much work a quality target demands; the \hl{Fleet Law}
says how long the fleet takes to deliver it, once communication, coordination and failure take their share.
\end{bigidea}

\begin{leftcol}
\sheetsec{Key ideas}
\begin{ideas}
\idea{Leaving one machine is a discontinuity.}{The binding limit moves from local memory to the
  network, then to power, and failure becomes routine.}
\idea{A synchronous step waits for its slowest worker.}{One straggler stalls every peer; a lost
  worker forces a choice: halt (consistent) or go on with stale updates (available).}
\idea{Scaling laws set the demand.}{Each halving of the loss above its floor costs the same resource
  multiple: a predictable return that keeps fleets growing.}
\idea{C\textsuperscript{3} splits step time.}{Compute (the math), Communication (the wire),
  Coordination (the logic). Relieve one and another binds.}
\idea{Failure is the steady state.}{A job that needs every device at once is rarely whole: design
  for fast, automatic recovery.}
\idea{The fleet stack places each decision.}{Infrastructure, distribution, serving, governance:
  each band constrains the ones above.}
\end{ideas}

\vspace{10pt}
\sheetsec{Which axis binds? Read the evidence}
\begin{sheettable}{l L L l}
\thead{Axis} & \thead{Evidence} & \thead{Relief} & \thead{Archetype} \\ \midrule
Compute & low share of peak, memory busy & fuse kernels & C (MobileNet) \\
Communication & links busy, slow collectives & fewer bytes, overlap & A (Llama-3) \\
Coordination & barrier waits, restarts & fix the slow node, recover & B (DLRM) \\
\end{sheettable}

\vspace{10pt}
\fcard{Scaling law}%
  {\ScalingLawLoss(\ScaledResource) = \ScaleCoefficient\,\ScaledResource^{-\ScalingExponent} + \IrreducibleLoss}%
  {Loss falls with resource $\ScaledResource$ (compute, data, parameters) toward a floor $\IrreducibleLoss$;
   halving the reducible part multiplies $\ScaledResource$ by $2^{1/\ScalingExponent}$. Training costs
   $\TrainingCompute \approx \val{flopsper}\,\Parameters\,\DatasetSize$ FLOPs for $\Parameters$ parameters and
   $\DatasetSize$ tokens; loss is lowest near $\DatasetSize \approx \val{tokensper}\,\Parameters$.}%
  {Too little data: data-limited, so buy data, not accelerators.}{}
\end{leftcol}%
\begin{rightcol}
\sheetsec{Formulas you will use}
\fcard{Fleet Law}%
  {\DistributedStepTime(\DeviceCount) = \under{\frac{\IronLawComputeTime}{\DeviceCount}}{compute}
   + \under{\CommunicationTime(\DeviceCount)}{communication} + \under{\SynchronizationTime(\DeviceCount)}{coordination}
   - \under{\OverlappedTime}{hidden}}%
  {Adding devices divides only the compute term; exchanges and waits stay, and less of them can hide.}{}{}
\fcard{How much of the step is useful}%
  {\UsefulComputeFraction = \frac{\IronLawComputeTime/\DeviceCount}{\DistributedStepTime(\DeviceCount)},
   \qquad \ScalingEfficiency = \frac{\SingleDeviceRunTime}{\DeviceCount\,\MultiDeviceRunTime}}%
  {Useful share of one step; speedup per device ($\SingleDeviceRunTime$, $\MultiDeviceRunTime$: one and
   $\DeviceCount$ devices). Name the measure you quote.}{}{}
\fcard{Reliability gap}%
  {\JobAvailability = \ComponentAvailability^{\,\DeviceCount}, \qquad
   \SystemMtbf = \frac{\ComponentMttf}{\DeviceCount}}%
  {All $\DeviceCount$ devices up at once; failures come $\DeviceCount$ times as often.}%
  {Independent failures. One more nine only delays the collapse.}%
  {\begin{tikzpicture}[sheet]
     \foreach \e/\t in {1/10,2/100,3/{1,000},4/{10,000}}{
       \draw[gridline] ({30+\e*50},0) -- ({30+\e*50},44);
       \node[small, anchor=north] at ({30+\e*50},-3) {\t};}
     \foreach \yv in {0,50,100}{\draw[gridline] (30,{\yv*0.44}) -- (230,{\yv*0.44});
       \node[small, anchor=east] at (27,{\yv*0.44}) {\yv\%};}
     \draw[muted, line width=0.9pt] plot[domain=0:4, samples=60, smooth] ({30+\x*50},{44*exp(10^\x*\val{lnfour})});
     \draw[accent, line width=1.1pt] plot[domain=0:4, samples=60, smooth] ({30+\x*50},{44*exp(10^\x*\val{lnthree})});
     \filldraw[fill=accent, draw=accent] (180,{44*\val{allupthree}/100}) circle (2pt);
     \filldraw[fill=white, draw=muted] (230,{44*\val{allupfour}/100}) circle (2pt);
     \node[lab, anchor=east] at (176,{44*\val{allupthree}/100-6}) {\val{threenines}\% per device};
     \node[note, anchor=east] at (190,32) {\val{fournines}\%};
     \node[note, anchor=north] at (130,-12) {devices required (log); y: all up};
   \end{tikzpicture}}
\fcard{Communication intensity}%
  {\CommunicationIntensity = \frac{\NetworkTrafficVolume}{\Operations}
   \quad\text{vs.}\quad \text{balance} = \frac{\LinkBandwidth}{\PeakRate}}%
  {Bytes sent per local operation, against link bandwidth over peak rate. Above balance the transfer
   ($\NetworkTrafficVolume/\LinkBandwidth$) outlasts the math ($\Operations/\PeakRate$).}%
  {Fewer bytes lower $\CommunicationIntensity$; overlap only hides the transfer.}%
  {\begin{tikzpicture}[sheet]
     \draw[dim] (10,0) -- (240,0);
     \node[small, anchor=north west] at (26,-4.5) {B/FLOP};
     \foreach \e/\t in {-6/{$10^{-6}$},-5/{$10^{-5}$},-4/{$10^{-4}$}}{
       \pgfmathsetmacro\x{10+(\e+6)*115}
       \draw[dim] (\x,0) -- (\x,-2.5);
       \node[small, anchor=north] at (\x,-4) {\t};}
     \pgfmathsetmacro\xs{10+(\val{cismalllog}+6)*115}
     \pgfmathsetmacro\xl{10+(\val{cilargelog}+6)*115}
     \pgfmathsetmacro\xb{10+(\val{balancelog}+6)*115}
     \draw[accent, dashed, line width=0.7pt] (\xb,-2) -- (\xb,24);
     \node[note, anchor=south, align=center] at (\xb,25) {balance\\ $\val{balance}$};
     \filldraw[fill=accentsoft, draw=accent, line width=0.5pt] (\xs,0) circle (2.4pt);
     \filldraw[fill=accent, draw=accent] (\xl,0) circle (2.4pt);
     \node[note, anchor=south, align=center] at (\xs,5) {\val{smallgpus} GPUs\\ $\val{cismall}$};
     \node[note, anchor=south east, align=center] at ({\xl+4},5) {\val{largegpus} GPUs\\ $\val{cilarge}$};
   \end{tikzpicture}}
\end{rightcol}

% ===================================================================== BACK: PRACTICE
\backside
\begin{answers}
\ans{Finish}{\val{halved} + \val{commexposed} = \val{afterstep} s; \val{smallstep} $\div$ \val{afterstep} =
  \val{afterspeedup}$\times$ faster, not \val{accel}$\times$; compute \val{newcompute} $\div$ \val{afterstep} = \val{afterfraction}\%.
  The exchanges do not speed up (Fleet Law).}
\ans{1}{(a) \val{qavail}\% to the power \val{qnodes} $\approx$ \val{qallup}\% chance all are up at once.
  (b) \val{mttf} h $\div$ \val{qnodes} = \val{qmtbf} h (reliability gap).}
\columnbreak
\ans{2}{$\val{qbytesfull} \div \val{qflop} = \val{qci}$ B/FLOP, above the balance \val{link} GB/s $\div$
  \val{peak} TFLOP/s $= \val{balance}$, so the link sets it (\val{qlink} s vs \val{qmath} s of math).
  Overlap hides time, not bytes. Half the bytes: $\val{qcihalf}$, below balance.}
\ans{3}{Coordination: wire unchanged, and the one worker that never waits is the straggler. Check its
  hardware and heat; replace or rebalance it (evidence table). Asynchronous updates drop the barrier wait
  but apply stale gradients: availability bought with consistency (straggler key idea).}
\ans{4}{(a) $\val{qpow}$ = \val{qmult}$\times$ the resource. (b) $\Parameters = \sqrt{\val{qbudget} \div \val{qkt}} \approx$ \val{qparams} billion parameters and $\DatasetSize \approx$ \val{qtokens} trillion tokens.
  (c) Data-limited: buy data (or train smaller), not accelerators (scaling-law card).}
\end{answers}

\begin{leftcol}
\sheetsec{Worked example}
\begin{practicebox}\boxtitle{Why do \val{gpugrowth}$\times$ the GPUs give less than \val{gpugrowth}$\times$ the speed?}
A modeled \val{params}B-parameter training step, global batch \val{batchseqs} sequences. On
\val{smallgpus} GPUs: compute \val{smallcompute} s, tensor-parallel exchange \val{smalltensor} s, exposed
gradient exchange \val{smallgrad} s, pipeline bubble (stages idle) \val{smallbubble} s. On \val{largegpus} GPUs, same batch:
\val{largecompute}, \val{largetensor}, \val{largegrad} and \val{largebubble} s (three decimals so the sum reproduces).\par\vspace{5pt}
{\centering
\begin{tikzpicture}[sheet]
  \pgfmathsetmacro\sc{165/\val{smallstep}}
  \pgfmathsetmacro\a{48+\val{smallcompute}*\sc} \pgfmathsetmacro\b{\a+(\val{smalltensor}+\val{smallgrad})*\sc}
  \pgfmathsetmacro\c{\b+\val{smallbubble}*\sc}
  \node[note, anchor=west] at (0,30) {\val{smallgpus} GPUs};
  \draw[accbox] (48,23) rectangle (\a,37); \node[lab] at ({(48+\a)/2},30) {compute};
  \draw[nbox] (\a,23) rectangle (\b,37);
  \draw[hotbox] (\b,23) rectangle (\c,37);
  \pgfmathsetmacro\p{48+\val{largecompute}*\sc} \pgfmathsetmacro\q{\p+(\val{largetensor}+\val{largegrad})*\sc}
  \pgfmathsetmacro\r{\q+\val{largebubble}*\sc}
  \node[note, anchor=west] at (0,8) {\val{largegpus} GPUs};
  \draw[accbox] (48,1) rectangle (\p,15);
  \draw[nbox] (\p,1) rectangle (\q,15);
  \draw[hotbox] (\q,1) rectangle (\r,15);
  \node[lab, anchor=west] at ({\r+5},8) {= \val{largestep} s; bubble is solid};
  \node[lab, anchor=south east] at (\c,38) {= \val{smallstep} s};
\end{tikzpicture}\par}\vspace{4pt}
\begin{steps}
\item Fleet Law: the exchanges are communication already net of overlap, the bubble is coordination
  (waiting). Sum: \res{\val{smallstep} s} and \res{\val{largestep} s}.
\item Speedup: \val{smallstep} $\div$ \val{largestep} = \res{\val{stepspeedup}$\times$} for \val{gpugrowth}$\times$ the GPUs.
\item Compute-time fraction: \val{smallcompute} $\div$ \val{smallstep} = \val{smallfraction}\%;
  \val{largecompute} $\div$ \val{largestep} = \res{\val{largefraction}\%}, here also the scaling efficiency
  (compute divides perfectly, so $\SingleDeviceRunTime/\DeviceCount = \IronLawComputeTime/\DeviceCount$).
  The unchanged bubble is now \val{largebubbleshare}\% of the step.
\end{steps}
\takeaway{More devices divide only the work that divides; the fixed terms set the speedup.}
\end{practicebox}



\vspace{10pt}
\sheetsec{Common mistakes}
\mistake{``Each GPU is reliable, so the job is reliable.''}%
  {One GPU fails every \val{mttf} h; a \val{fleetgpus}-GPU job, every \val{mttf} $\div$
   \val{fleetgpus} = \val{mtbf} h, about \val{perday} times a day.}{}
\mistake{``The compute budget predicts the training time.''}%
  {Scaling laws price work, not time. Spread a fixed run for \val{planned}$\times$ the speed; if communication
   and coordination take \val{overhead}\% of the scaled step, you get \val{planned} $\times$ \val{keepshare} = \val{realized}$\times$.}{}
\end{leftcol}%
\begin{rightcol}
\sheetsec{Your turn: finish it}
\begin{practicebox}\boxtitle{Accelerators twice as fast, on \val{smallgpus} GPUs}
New accelerators do arithmetic \val{accel}$\times$ faster; the exchanges
(\val{smalltensor} + \val{smallgrad} = \val{commexposed} s) do not change.\par\vspace{4pt}
\textbf{Given:} compute and bubble halve: \val{newcompute} + \val{halfbubble} = \val{halved}~s.\par
\textbf{You:} find the new step, the speedup and the new compute-time fraction, and say why it is not \val{accel}$\times$.\par\vspace{30pt}
\end{practicebox}

\vspace{10pt}
\sheetsec{Check yourself}
\begin{qs}
\item \assumed{} a job needs \val{qnodes} nodes, each \val{qavail}\% available with an MTTF of
  \val{mttf} h. (a) Chance all are up at once? (b) Time between job failures?
\item \assumed{} per step a worker sends \val{qbytes} GB and does $\val{qflop}$ FLOPs; link \val{link} GB/s,
  peak \val{peak} TFLOP/s. Does the link or the math set the step? Does overlap lower the ratio? Do
  8-bit gradients (from 16-bit)?
\item A cluster doubles and steps slow down. Links and collectives look unchanged, but barrier waits
  grew on every worker but one. Which axis binds? What next? Would going asynchronous remove the wait, and at what cost?
\item \assumed{} exponent \val{qalpha}. (a) What multiple of resource halves the reducible loss?
  (b) Compute-optimal model and data for $\val{qbudget}$ FLOPs? (c) Only \val{qhave} billion tokens
  exist: which regime, and what do you buy?
\end{qs}
\end{rightcol}
\end{document}
